% LAB BOREDOM — module L1, The Three Literatures. GUIDED DRAFTING (desktop recipe), Section 2 run.
% Session 2026-08-06. Author rulings: pack v2 is the bank; r1 ¶1 carried as working base.
% Pack: lab-boredom-L1-pack-v2 (+ -v2-quotes.json)  |  Spec: outline v7 §L1  |  Rulings: register v8.
% Protocol: Stage-3 guided batches (1–3 ¶¶, approval gate each batch); «Qnn» markers, mechanical
%   substitution; rendered copy = L1-DRAFT-v1-rendered.tex. MASTER FILE — edit this, re-render after.
%
% STANDING RULINGS ADDED THIS SESSION (author notes on batch 1):
%   (i) voice must match Part III Section 1 (desktop register, its cadence);
%   (ii) FIRST + LAST NAME at an author's first mention in the section, surname after;
%   (ii-b) ★ STANDING (2026-08-06): DO NOT USE "instrument"/"instruments" AS A VAGUE COVER.
%        Name the thing: the eye tracker, the scalp electrodes, gaze deviation, eye closure,
%        pupil diameter, the seven-item survey, "the physiological measurements." The word is
%        permissible only where the source itself uses it (Holmqvist's reporting factors) or
%        where the referent was named in the same breath. AUDIT THIS BEFORE EVERY BATCH.
%        Related: EEG and eye trackers are general-purpose recording equipment — WE DID NOT
%        BUILD THEM FOR BOREDOM. Never write "instruments built to detect boredom." The
%        correct form: physiological measurements scored against self-reports of boredom
%        collected under Eastwood's definition.
%   (ii-e) ★ STANDING (author, 2026-08-07): REFER TO A SPECIFIC SUBJECT BY NUMBER, NOT BY
%        PRONOUN. Write "S04 fought sleep at six against one", never "they fought sleep".
%        Pronouns for a named subject are avoided entirely so the reader never has to track an
%        antecedent across sentences. Everything is anonymized, so the number IS the name.
%   (ii-d) ★★ STANDING RULE (author, 2026-08-07) — COMPLETE-COHORT ANALYSIS BEFORE ANY
%        RESULTS CLAIM. Never draft a claim about results from one participant unless the
%        passage is explicitly about that participant. Run EVERY subject, EVERY time, and
%        review the full breakdown before writing. Origin: the S04 case paragraph was drafted
%        from S04 alone and badly understated the cohort — the full table showed each measure
%        agreeing with participants own reports in 10 or 11 of 12 cases, and showed that two of
%        the three apparent exceptions dissolve on inspection.
%   (ii-c) ★ PUPIL DEMOTED IN L1 BY AUTHOR RULING 2026-08-06: pupil dilation yields nothing
%        L1's argument builds on. It does not track self-report (-0.172 vs boredom, ns), and
%        its BOR-vs-CLC separation admits several uncontrolled explanations at once (luminance
%        unmatched; arousal; effort). Measured means: CLC 3.51 mm > INT 3.40 > BOR 3.32 — the
%        OPPOSITE of what a pure brightness account predicts if clinical is the brightest
%        sample, which is itself the evidence that no single factor governs it.
%        ★ TENSION TO RESOLVE AT L3: register v8 R-05 lists pupil among the LOAD-BEARING
%        channels, and L3.1 records that nine Holm-surviving contrasts fall on eye closure,
%        pupil, or excursion frequency. L1 now says the argument does not rest on pupil.
%        These are compatible only if L3 reports pupil's separation WITHOUT building on it.
%        Decide explicitly at L3 rather than letting the two modules drift apart.
%   (iii) pack first, but the rebuilt archon-cli-v3 index is an encouraged evidence surface —
%        used this revision for: self-citing author (citation-network §2), DA-03 roster (debate-map),
%        Quaranta claims (bor-sec-10), regulatory-paper theses (bor-sec-25/26), P1/P2 citations
%        (king-salvo-physio-2023 + grp-program-lineage), Freeman/Withy/McManus first names.
%
% BATCH 2 = ¶¶4–5 (drafted 2026-08-06, awaiting author review).
%   ¶4 (education): Mansikka three-way typology 262–66, second form + das Man 264 (bor-sec-18);
%     Gibbs datable-time/concealment 601–02 (bor-sec-07); Feldges & Pieczenko + Mertel +
%     Aroles & Küpers compactly in a footnote (bor-sec-16 and units); THOMSON GUARD APPLIED
%     per A-02 — cited for paideia/enframing/letting-learn frame, stated NOT a boredom source.
%   ¶5 (EEG): Kim/Seo/Laine null at 13 participants (A-06, bor-sec-28) — framed as the study
%     that sought agreement and documented disagreement; Barry et al. baselines 2765–66,
%     2771–72 (A-07, bor-sec-27); Yuvaraj et al. «Q16» at printed p. 11 (A-08, bor-sec-33)
%     with the no-resting-baseline gap; Miyauchi & Kawasaki, Perone, Yakobi, Seo compactly.
%   Names verified from index citations: Jan-Erik Mansikka, Paul Gibbs, Tom Feldges + Sonia
%     Pieczenko, Kevin Mertel [FIRST NAME ****** — index has "Mertel, K." only], Jeremy Aroles,
%     Wendelin Küpers, Iain Thomson, Joochan Kim, Jungryul Seo, Teemu H. Laine, Robert J. Barry,
%     Rajamanickam Yuvaraj, Eri Miyauchi, Masahiro Kawasaki, Sammy Perone, Ofir Yakobi.
%   Q16 is the batch's only rendered quotation. A-10 respected: no Miyauchi quotation exists.
%   Kim's null PARAPHRASED from the unit digest (no gated verbatim in bank) — E05 CONFIRMED.
% BATCH 2 REV 1 (2026-08-06, author notes):
%   ¶4a: the two-channel point EXPANDED — names the actual apparatus (continuous eye tracking +
%     EEG through the film; 7-item self-report after each episode per R-07), states Mansikka's
%     evidentiary limit, and lands the compatibility claim as a measured divergence on the
%     clinical footage (keystone structure per pack P6; NO numbers quoted here — L3/L4 carry them).
%   ¶4b: Thomson correction MOVED into the compact-citations footnote (author: not pertinent to
%     the body's argument). A-02 guard preserved verbatim in the footnote.
%   ¶5a: "instrumented literature" -> the field NAMED (psychophysiology and affective computing,
%     per the unit disciplines: Kim biomedical signal processing/affective computing; Barry
%     clinical neurophysiology; Yuvaraj cognitive neuroscience/educational psychology).
%   ¶5b: "arrives carrying its nulls" REPLACED with the concrete result (Kim's hypothesis and
%     its failure, stated before it is characterized).
%   ¶5c: CHANNEL now GLOSSED at first use in this module (pack P6's locked definition rendered
%     in prose) so later modules can use the term as defined.
%   ¶5d: "fill out the strand" REPLACED — each of the four compact studies now states its own
%     contribution (Miyauchi key-press onset marking; Perone FEA as dynamic regulation not trait;
%     Yakobi state/trait dissociation; Seo 86.73% k-NN on two electrodes -> rendered "86.7").
%
% BATCH 2 REV 2 (2026-08-06, author notes): Thomson guard sentence DELETED from fn d (frame
%   attribution retained); "(EEG)" added at first use of electroencephalographic; the
%   four-compact-studies sentence MOVED to a footnote.
%
% BATCH 3 = ¶¶6–7 (drafted 2026-08-06, awaiting author review).
%   ¶6 (eye tracking): Holmqvist et al. 2023 minimal reporting guideline, 47 authors, 207-study
%     comparison (bor-sec-36) — L2 is measured against it; van den Brink/Murphy/Nieuwenhuis
%     pupil-tracks-lapses + inverted-U adaptive gain (D-18 entry) = what pupil indexes;
%     SCHARINGER A-04 oppositional-ranking precedent with r = -.16/-.14 ns, alpha alone tracking
%     comprehension; YEAR CORRECTION 2015-not-2019 carried in a footnote (A-03); VanderWerf
%     gaze-position dependency + «Q19» at printed 2785 + unknown-mixture caution (E13);
%     Charoenpit, Jaques, Sharma, GazeMotive, Cheval, Skaramagkas review in a compact footnote.
%   ¶7 (negative through-line): Kim null; Perone's FIRST hypothesis unsupported (distinct from
%     the batch-2 footnote's second-hypothesis fact); Yakobi state/trait dissociation + «Q17»;
%     Lindley & Nacke GEQ 2-of-7 with flow/immersion failing to separate + «Q21» in footnote;
%     Miyauchi appraisal-not-boredom; closing = this section's EEG negative is characteristic.
%
% BATCH 4 = ¶¶8–9 territory, ¶8 split into TWO paragraphs (drafted 2026-08-06, awaiting review).
%   Spec: outline v7 §L1 ¶8. O-06 SETTLED = cite rather than re-present, so neither published
%   study is summarized at length; both are characterized and their two moves named.
%   ¶8a (the lineage): Group B `grp-program-lineage` — trajectory 2024 applicability/satisfaction
%     -> 2023 physiological -> 2025 scaling; warrant identical across five years (hospital space,
%     safety protocols, cohort limits = scarcity of clinical access); published position =
%     viability settled + experience favorable WHILE the attention/engagement question stays
%     open in the program's own record (E15). Extension rather than resolution.
%   ¶8b (the two moves): definitional adoption Eastwood -> Fahlman -> P1 -> this data (E16/P6
%     spine; Q05's content PARAPHRASED here since ¶1 already carries the adoption — the verbatim
%     definition is held for ¶10 per D-23); then the redemptive half — «Q30» + refs [13-15]
%     rendered as VR-PHEN records it (E14: one undifferentiated gesture, invoked but inert).
%   D-20 applied: the one-of-fifty-five observation (E18 neighborhood) stated in a footnote as
%     DESCRIPTION not indictment, explicitly labeled as such.
%   First names verified from Group A entry: Jacquelyn Morie; Heinzel & Heinzel (Sibylle and
%     Stephan — ****** VERIFY at assembly); Jason Tham (****** VERIFY).
%   NEW BANK ENTRY Q30 (P1's framing sentence) — verified EXACT this session, and verified to
%     sit at the END OF THE INTRODUCTION (before Methods), so the no-abstract lock is satisfied.
%   P1 stays quoted rather than characterized (E16) — Group B deliberately excludes it.
%
% BATCH 3 REV 4 (2026-08-06, author challenge — CHECKED AGAINST THE v5 TREE, author correct):
%   "its measures lean the same way as every other measure in the study" was WRONG twice over.
%   Checked in boredom-analysis-v5 FINDINGS.md + APPENDIX-CRITERION-VALIDITY.md:
%   (a) The eye measures DID reach significance. Criterion validity, n=36, permutation p:
%       gaze deviation (per-participant) +0.590 (0.0024) vs boredom, -0.585 (0.0022) vs
%       engagement; per-file +0.524 (0.0074) / -0.527 (0.0075); eye closure +0.460 (0.0212) /
%       -0.416 (0.0370). So "every other measure" only leaning is false.
%   (b) EEG IS ABSENT FROM THE CRITERION TABLE ENTIRELY — it is R1-only N=8 while the criterion
%       analysis pools 36 episodes across rounds, so EEG cannot enter it. Any sentence
%       comparing EEG's agreement to the other channels' agreement is comparing to nothing.
%   EEG's actual best showing (FINDINGS.md): DMN power (alpha+theta), Friedman p=0.1969;
%       BORvCLC and BORvINT both Wilcoxon 0.078 / Holm 0.2344, rb +0.722, direction 6/8;
%       CLCvINT 0.5469. Frontal alpha p=0.4169, frontal theta p=0.8825, parietal alpha p=0.2231.
%   Prose now: eye measures separated AND agreed (findings); EEG gave a 6-of-8 direction on the
%   boring-film contrasts = a possible relationship, reported as one, not a result; plus the
%   R1-only thinness. NO per-channel significance claimed (register v8 standing rule).
%
% BATCH 3 REV 3 (2026-08-06, author notes — TWO SUBSTANTIVE CORRECTIONS):
%   ¶6: CLINICAL FOOTAGE DE-EMPHASIZED. The divergence example naming the clinical film is
%     removed; the point is now stated for any film. AUTHOR RULING: this section analyzes
%     boredom vs engagement writ large with all three samples treated equally — the clinical
%     stimulus does NOT hold the prominent position it held in the published study. Conclusions
%     about the clinical film remain permissible; foregrounding it does not.
%     ★ CARRY THIS INTO L3/L4 — it bears on how the keystone is presented.
%   ¶7: "The EEG half of that design returned nothing" WAS FACTUALLY WRONG (author). Corrected:
%     the EEG leans the same way as every other measure — a possible relationship worth
%     reporting — but reaches significance nowhere. Grounded in the O-9 reprocessing at N=8
%     (DMN p=0.20, parietal alpha p=0.22; report as convergent DIRECTION, never as a
%     per-channel result) plus outline v7 §L3.8 (0 of 5 surfaces; 0 of 15 contrasts through
%     Holm). Permitted form now used: "a direction, not a finding."
%     The 4-of-8 concordance figure is left to L3 rather than stated here.
%
% BATCH 3 REV 1 (2026-08-06, author notes):
%   ¶6a: "inverted U" REPLACED with a plain-language statement of the curve (best at
%     intermediate diameter; both extremes degrade attention).
%   ¶6b: fn e reduced to the citation + why it is cited; the year is simply given as 2015.
%   ¶6c: «Q19» REMOVED from prose. Reason: its span begins with a capital T and runs as a
%     sentence fragment, so the bracket-lowercase rule requires "[t]o...", which mechanical
%     substitution cannot produce. Per the pack's own recommendation for VanderWerf
%     (report, do not quote), the paragraph now REPORTS the blink measurement (334 ± 67 ms,
%     Q18's content, rendered with \pm not U+2AFE) and EXPLAINS the attention control in
%     prose. Locus widened to 2785--86 (Q19 at 2785, Q18 at 2786).
%   ¶6d: GAZE literature MOVED INTO THE BODY at author's direction — Jaques, GazeMotive
%     (Wang/Xu/Chen), Sharma/Giannakos/Dillenbourg, Charoenpit & Ohkura, Cheval, Skaramagkas
%     review; compact footnote deleted. The added analytic point: all are classifier designs
%     validated AGAINST self-report, so the report is their arbiter — the arrangement this
%     study cannot adopt, since its question is what happens when eyes and report disagree.
%   ¶7a: the EEG negative STATED SPECIFICALLY (0 of 5 surfaces, 0 of 15 contrasts through
%     Holm, concordance 4 of 8 = chance) from outline v7 §L3.8. Numbers only; L3 carries detail.
%   ¶7b: STATE/TRAIT now DEFINED at first load-bearing use, in a footnote, with «Q29»
%     (Raffaelli 2452, verified EXACT this session). ¶5's footnote reworded to plain language
%     so the technical terms are defined where they carry weight.
%   ¶7c: fn g stripped of the metatextual correction note — it now cites Lindley and Nacke
%     correctly and says nothing else. The correction is recorded HERE and in the errata file.
%   ¶7d: EEG-DROP added per author — first person, our own decision, noise + scalability,
%     with «Q26»/«Q27»/«Q28» from P2's INTRODUCTION (not the abstract), all verified EXACT
%     this session. D-16 rendered as: dropped for scaling, before the reprocessing established
%     what the measurements could have carried.
%
% ★ SESSION BANK: substitution now consumes `L1-bank-session.json` = pack v2 bank + the
%   four in-session additions in `L1-addendum-quotes.json` (Q26-Q29). Pack v2 remains
%   untouched; fold the addendum into a pack v3 before another module reuses Q26-Q29.
%
% ★ CITATION-FORM CORRECTION FOUND THIS BATCH (index-verified, pack is wrong):
%   Q21's cite string reads "Csikszentmihalyi, qtd. in Nacke and Craig" and the corpus directory
%   is named "Nacke & Craig". CRAIG IS A GIVEN NAME. The index citation field reads
%   "Lindley, Craig A., and Lennart Nacke" — bor-sec-44. Prose renders LINDLEY AND NACKE and the
%   error is stated in a footnote. SAME CLASS OF TRAP as the Colin/Barrett case. Pack v2 errata
%   item; fix the cite string before any later module reuses Q21.
%   NB Q21 rendered with "(qtd. in Lindley and Nacke 1--2)" — the bare «Q21» form is used, NOT
%   «Q21+», because the bank's cite string carries the wrong surname.
%
% BATCH 1 = ¶¶1–3. REV 1 (2026-08-06, per author notes):
%   ¶1: Elpidorou NAMED as the self-citing author (citation-network.md §2: 7 of 8 A→C links are
%       Elpidorou citing Elpidorou; 8th = Hughes's critical footnote) + fn with both regulatory works
%       summarized; first "we" now names Christine King + fn listing P1/P2 and the appendix; ending
%       rebuilt around the Fahlman→Eastwood adoption and the design it governed.
%   ¶2: ****** resolved into DA-03 footnote (O'Brien/Hadjioannou/McManus + SZ 310 vs Withy + FCM
%       143/165; E&F periodizing at Affectivity I 662n2); Slaby point expanded via the disposedness
%       (Befindlichkeit) argument (link to Part I rendered through shared vocabulary, not explicit
%       back-reference; BT 172–75 locus carried from Part I draft); first-mention names applied.
%   ¶3: "VR headset"; conclusion rebuilt on Quaranta's actual apparatus (Loht's viewing-Dasein,
%       Drive My Car dead-time aesthetics, mood as neither work-alone nor viewer-alone, 31–32, 38).
%
% Batch-1 verification ledger:
%   Q14 printed p. 117 PINNED (folio-read: running head "The Other Side of Existence 117"; Slaby
%     fn 18 = Pessoa, Book of Disquiet, fragment 219). Q07 vol. pp. 207–08; Q12 printed p. 101.
%   P1 full citation from king-salvo-physio-2023-00-overview (ASEE 2023 #37129, 5 authors).
%   P2 from VLE manifest: "King et al., Assessment of Student Engagement in VR Clinical Immersion
%     Environments through Eye Tracking, ASEE 2024 #44685" — FULL AUTHOR LIST PENDING at assembly.
%   Elpidorou regulatory theses from bor-sec-25/26 digests (PhenomCogSci 17: 455–84; PhilPsych
%     31.3: 323–51). DA-03 roster from debate-map.md. Quaranta from bor-sec-10 (Film-Philosophy
%     28.1: 31–46); "neither subjective nor objective" p. 38 PARAPHRASED (quote not gated).
% PENDING VERIFICATION (assembly): first names for O'Brien, Hadjioannou (ed./ch. author), Hughes;
%   McManus chapter title as printed; Slaby 2010 venue for WC; SZ-vs-BT page convention harmonize.
% OPEN: Appendix designation for P1/P2 reprints (****** in fn A). E01 open item CLOSED this rev
%   (author = Elpidorou, index-verified).
% Selection decisions carried: Q01 barred; Q08 held for ¶10 (paraphrased here); Q09 paraphrased;
%   Q10 unused (optional); Q11 skipped; Q13 avoided (hazard); O-06 = cite-not-re-present.
%
% WORKS CITED introduced by this batch (for assembly; ****** = verify there):
%   Eastwood, John D., et al. "The Unengaged Mind: Defining Boredom in Terms of Attention."
%     Perspectives on Psychological Science, 2012. [vol/no/pp ****** at assembly]
%   Elpidorou, Andreas. "The Bored Mind Is a Guiding Mind: Toward a Regulatory Theory of Boredom."
%     Phenomenology and the Cognitive Sciences, vol. 17, 2018, pp. 455-84.
%   ---. "The Good of Boredom." Philosophical Psychology, vol. 31, no. 3, 2018, pp. 323-51.
%   Elpidorou, Andreas, and Lauren Freeman. "Affectivity in Heidegger I: Moods and Emotions in
%     Being and Time." Philosophy Compass, 2015. [vol/pp ******]
%   ---. "Is Profound Boredom Boredom?" Heidegger on Affect, ed. Hadjioannou, 2019, pp. 177-[**].
%   Fahlman, Shelley, et al. "Development and Validation of the Multidimensional State Boredom
%     Scale." 2013. [journal details ******]
%   Hughes [first name ******]. "Meaninglessness and Monotony in Pandemic Boredom." 2023.
%   King, Christine E., Kit Roy Feeney, Quangminh Tang, Milan Das, and Dalton Salvo. "Work in
%     Progress: Physiological Assessment of Learning in a Virtual Reality Clinical Immersion
%     Environment." ASEE, 2023. Paper #37129.
%   King, Christine E., et al. [****** full list] "Assessment of Student Engagement in VR Clinical
%     Immersion Environments through Eye Tracking." ASEE, 2024. Paper #44685.
%   O'Brien [******]. "Being, Nothingness and Anxiety." Heidegger on Affect, 2019.
%   Hadjioannou [******], ed. Heidegger on Affect. Rowman?, 2019. [publisher ******] — also ch. 4,
%     "Angst as Evidence: Shifting Phenomenology's Measure."
%   McManus, Denis. [ch. 6 title per index: "Affect and Authenticity"] Heidegger on Affect, 2019.
%   Withy, Katherine. "Finding Oneself, Called." Heidegger on Affect, 2019.
%   Quaranta, Chiara. "In the Mood for Heideggerian Boredom? Film Viewership as
%     Being-in-the-World." Film-Philosophy, vol. 28, no. 1, 2024, pp. 31-46.
%   Slaby, Jan. "The Other Side of Existence: Heidegger on Boredom." 2010, pp. 101-20. [venue ******]
% --- introduced by batch 2 ---
%   Aroles, Jeremy, and Wendelin Küpers. "Towards an Integral Pedagogy in the Age of 'Digital
%     Gestell'." 2022. [venue/pp ******]
%   Barry, Robert J., Adam R. Clarke, Stuart J. Johnstone, Christopher A. Magee, and Jacqueline
%     A. Rushby. "EEG Differences Between Eyes-Closed and Eyes-Open Resting Conditions."
%     Clinical Neurophysiology, vol. 118, 2007, pp. 2765-73.
%   Feldges, Tom, and Sonia Pieczenko. "Boredom in Educational Contexts: A Critical Review."
%     Encyclopaideia, vol. 24, no. 57, 2020, pp. 1-15.
%   Gibbs, Paul. "The Concept of Profound Boredom: Learning from Moments of Vision." Studies in
%     Philosophy and Education, vol. 30, 2011, pp. 601-13.
%   Kim, Joochan, Jungryul Seo, and Teemu H. Laine. "Detecting Boredom from Eye Gaze and EEG."
%     Biomedical Signal Processing and Control, vol. 46, 2018, pp. 302-13.
%   Mansikka, Jan-Erik. "Can Boredom Educate Us? Tracing a Mood in Heidegger's Fundamental
%     Ontology from an Educational Point of View." Studies in Philosophy and Education, vol. 28,
%     2009, pp. 255-68.  [2008/2009 — confirm year at assembly ******]
%   Mertel, K[****** first name]. "Heidegger, Technology and Education." Journal of Philosophy
%     of Education, vol. 54, no. 2, 2020, pp. 467-86.
%   Miyauchi, Eri, and Masahiro Kawasaki. "Electroencephalogram Oscillations Support the
%     Involvement of Task-Unrelated Thoughts in the Mechanism of Boredom: A Pilot Study." 2018.
%     [venue/pp ******]
%   Perone, Sammy, Elizabeth H. Weybright, and Alana J. Anderson. "Over and Over Again: Changes
%     in Frontal EEG Asymmetry Across a Boring Task." 2019. [venue/pp ******]
%   Seo, Jungryul, Teemu H. Laine, and Kyung-Ah Sohn. "Machine Learning Approaches for Boredom
%     Classification Using EEG." 2019. [venue/pp ******]
%   Thomson, Iain. "Heidegger on Ontological Education, or: How We Become What We Are."
%     Inquiry, vol. 44, 2001. [pp ******]
%   Yakobi, Ofir, James Boylan, and James Danckert. "Behavioral and Electroencephalographic
%     Evidence for Reduced Attentional Control and Performance Monitoring in Boredom." 2021.
%     [venue/pp ******]
%   Yuvaraj, Rajamanickam, et al. "EEG-Based Functional Connectivity Patterns During Boredom in
%     an Educational Context." Scientific Reports, vol. 15, 2025, article 35439.

Boredom has been academically pursued across two rather different fields of study. One asks
what boredom discloses about the person who suffers it, treats the experience as a disclosure
of how things stand for someone rather than as a state to be relieved, and reaches for Martin
Heidegger and the phenomenological tradition that follows him; the other asks what boredom
looks like in a signal, and reaches for electrodes, eye trackers, and self-report scales
validated against one another. Across the works surveyed here, the connection between them
amounts to one philosopher, Andreas Elpidorou, citing his own work across his own two research
programs, together with one critical footnote written by someone else.\footnote{Elpidorou
publishes in both literatures, and the phenomenological essays---several written with Lauren
Freeman---cite two papers of Elpidorou's own regulatory-theory program. ``The Bored Mind Is a
Guiding Mind: Toward a Regulatory Theory of Boredom'' (\textit{Phenomenology and the Cognitive
Sciences}, 2018) argues that boredom is a functional emotion, at once informing a person that
the present situation has stopped satisfying and motivating the pursuit of a new goal. ``The
Good of Boredom'' (\textit{Philosophical Psychology}, 2018) argues that transitory state
boredom serves well-being precisely by keeping a person in motion toward goals worth having.
The one bridge not of Elpidorou's making runs the opposite direction and is a critique:
Emily Hughes, ``Meaninglessness and Monotony in Pandemic Boredom'' (2023), objects in a footnote
that Elpidorou's normative-psychological turn departs from Heidegger's project.} The first
literature counts a reading of experience
as evidence; the second counts a number extracted from a sensor. Neither treats what the
other counts as bearing on its own question, and the absence is not an oversight either could
have corrected by wider reading. The two studies Christine King and I published stand in
neither literature:\footnote{The two studies: King, Christine E., Kit Roy Feeney, Quangminh
Tang, Milan Das, and Dalton Salvo, ``Work in Progress: Physiological Assessment of Learning
in a Virtual Reality Clinical Immersion Environment'' (ASEE 2023, Paper \#37129), the
three-participant pilot that first instrumented the platform's videos---EEG, cardiac, and
gaze---against boring and interesting control films; and King, Christine E., Matthew Lo, Milan Das, and
Dalton Salvo, ``Assessment of Student Engagement in Virtual Reality Clinical Immersion
Environments through Eye Tracking'' (ASEE 2024, Paper \#44685), the eye-tracking study of
student engagement on the same platform. Both papers are reproduced in
full in Appendix ******.} they belong to a third lineage, clinical immersion in engineering
education, where a headset substitutes for clinical access that no program can secure in
sufficient quantity, and that lineage arrives with its own question, its own warrant, its own
reference list, and no commitment about what boredom is. From the second literature we took
one thing, deliberately and by name: the definition of boredom under which we collected the
data. Our first paper adopts Shelley Fahlman's formulation of it---the definition
underwriting the Multidimensional State Boredom Scale, which carried John Eastwood's
attentional account of boredom into measurement---and the adoption governed the design end to
end.\footnote{Fahlman and colleagues state it thus: ``«Q31»'' (69). They offer it
transtheoretically, as the common core of the existing theories rather than as one theory
among them, and that framing is what licensed its adoption elsewhere as a neutral starting
point. Eastwood is a co-author, and the attentional account arrives with the definition.} Because the definition fixes boredom as an aversive episode, an unfulfilled desire for
engagement that would satisfy, the experiment treats boredom as something a film can induce
and relieve: it sets a film selected to starve engagement beside a film selected to feed it
and beside the clinical footage the platform exists to deliver, and it points its
instruments---EEG, eye tracking, and a per-episode self-report---at what each film does to
the person watching. What the definition does not contemplate is a boredom that is not about
the film at all. The data collected under it contains exactly that.

% ****** AUTHOR ACTION PENDING (¶1 footnote b) — rendered so it cannot be lost in a compile.
% ****** Remove this box once the appendix exists and the footnote's ****** is filled.
\begin{center}
\fbox{\parbox{0.9\linewidth}{\centering\textbf{****** AUTHOR ACTION PENDING ******}\\
Add the appendix containing the two published articles (P1, ASEE 2023 \#37129; P2, ASEE 2024
\#44685), then fill the appendix designation in the footnote above. Also pending: P2's full
author list for the works cited.}}
\end{center}

The phenomenological literature's liveliest quarrel concerns whether the experience Heidegger
% [ERRATUM APPLIED 2026-08-15, author-approved: «Q07» cite corrected (207--08) -> (178).
%  The old pin was the verify tool's PDF-page field; folio-read = volume p. 178 (running
%  head "178  A. Elpidorou and L. Freeman"). Chapter folios are 177-203-region, NOT 207-233;
%  the standalone E&F PDF is an UNCORRECTED PROOF paginated 1-27. WC stub updated to
%  pp. 177-[**] (end folio verified at assembly). Backup: .backups/L1-erratum2-2026-08-15/]
analyzes is boredom at all. Elpidorou and Lauren Freeman set profound boredom against the
construct the empirical programs measure---a brief, situation-bound, many-sided response to
unsatisfying circumstances---and conclude of profound boredom that ``«Q07»'' (178). The
verdict protects both parties to the comparison: it spares the psychological construct a
metaphysics it never claimed, and it cautions every instrumented study, this one included,
against announcing that a sensor has caught what Heidegger meant. Nor is it settled which
mood holds the ground position in Heidegger's own architecture: the case for anxiety
(\textit{Angst}) as the fundamental attunement (\textit{Grundstimmung}) rests on
\textit{Being and Time}'s near-unique designation, the case for profound boredom on
\textit{The Fundamental Concepts of Metaphysics}, and Elpidorou and Freeman observe that the
question is partly one of period, since the term appears only once in \textit{Being and Time}
and becomes load-bearing only after 1929.\footnote{The anxiety side is argued across three
chapters of \textit{Heidegger on Affect} (ed. Christos Hadjioannou, 2019): Mahon O'Brien,
``Being, Nothingness and Anxiety''; Hadjioannou's own ``Angst as Evidence: Shifting
Phenomenology's Measure''; and Denis McManus on affect and authenticity---with \textit{Being and Time}
itself giving anxiety the designation nearly alone (SZ 310). Katherine Withy's chapter in the
same volume, ``Finding Oneself, Called,'' makes the affirmative case that profound boredom
can occupy the role, with support at \textit{FCM} 143 and 165. For the periodizing
observation: Elpidorou and Freeman, ``Affectivity in Heidegger I: Moods and Emotions in
\textit{Being and Time}'' (2015), 662n2.} Jan Slaby reads the whole affective register as
``«Q12»'' (101): mood is not noise laid over thinking but the way a person's standing in the
world shows itself at all. Slaby's formula names what \textit{Being and Time} calls
disposedness (\textit{Befindlichkeit})---one does not first register a neutral scene and
afterward color it with feeling; one finds oneself already attuned, affected by what one
encounters before any reflective act of judgment arrives (\textit{BT} 172--75)---and the
formula carries a consequence. If attunement is disclosure rather than coloration, then a
person can remain occupied---attentive, compliant, even performing well---while what the
attunement discloses is emptiness. Fernando Pessoa's words, as Slaby borrows them, give that
depth its starkest formula---``«Q14»'' (qtd. in Slaby 117). An occupied surface over a hollow
depth is, under this reading, a structure open to description, not a contradiction in terms.

Film is the point where the phenomenological literature reaches this experiment's own medium.
Chiara Quaranta, building on Shawn Loht's account of film viewing as a genuine mode of
being-in-the-world---a two-way disclosure in which the viewer opens onto the film and the
film discloses to the viewer---argues that Heidegger's three forms of boredom, the profound
form above all, name what certain films do to the people watching them. The case study is
Ryusuke Hamaguchi's \textit{Drive My Car}, whose fixed shots, silences, repetitions, and dead
time Quaranta reads as an aesthetics composed to arouse profound boredom---with the standing
qualification that no film accomplishes this alone, since a mood arises neither from the work
by itself nor from the viewer by itself but between the two (31--32, 37). Within the works
surveyed here this is the one sustained reading of film through the boredom analysis, and the
arrangement it describes is the arrangement this experiment builds and instruments: a film
selected for monotony, watched under mandate rather than by choice, in a VR headset that
holds the footage before the eyes, while EEG, eye tracking, and a per-episode self-report
record what the arrangement does to the person inside it. Quaranta reads the arrangement off
a finished film; the experiment assembles it on purpose and asks every channel to report.

Education is where this analysis has been put to work longest, and where its most useful
precedent for the present study is found. Jan-Erik Mansikka fits Heidegger's three forms to
three ways of learning: the first to technical instruction that reduces phenomena to
abstractions and leaves the student confronting dull matter; the second to a learning that is
meaningful and still leaves the student empty; the third to boredom as a possible rupture,
which Mansikka is careful not to prescribe as method (262--66). The second is the one that
matters here. A student, on his account, can be outwardly exemplary---fully occupied,
participating, even high-achieving---while what governs the performance is the anonymous
``they'' (\textit{das Man}) rather than any self-relation to the material (264). The
diagnostic consequence follows directly: engagement and this hollow boredom are compatible
rather than opposed, so a metric that counts engagement---time on task, submissions,
participation---cannot detect the second form, because everything such a metric registers is
precisely what the second form leaves intact. Mansikka argues this in a classroom register
more than a decade before this experiment, and argues it from the phenomenology alone: he has
the student's outward performance and Heidegger's analysis, and nothing that would let him
exhibit the emptiness underneath. The experiment's data supply what his argument lacks. Each
participant here watched under two kinds of observation at once---eye tracking and EEG running
continuously through the film, and a seven-item self-report administered after each
episode---so every episode yields both a measurement of the person's outward comportment and
that same person's account of what watching was like. Mansikka's compatibility claim predicts
that these two can come apart: a viewer whose eyes and body register occupation while the
account given afterward reports emptiness. That is what these recordings contain, and because
both sides of the divergence are captured on the same episode, the compatibility of engagement
with hollow boredom stops being an interpretation of a student's behavior and becomes a
measured result. Paul Gibbs pursues the same third form toward pedagogy, reading modern datable time---the
semester, the credit hour, the examination date---as the very thing that conceals boredom's
disclosive work, so that busyness is not boredom's opposite but the mechanism of its
concealment (601--02). Tom Feldges and Sonia Pieczenko press a different objection, that the
field's dominant taxonomy of academic boredom conflates what precedes boredom with boredom
itself, and that classroom boredom may be an artifact of frontal instruction rather than
anything intrinsic to learning; Kurt Mertel, Jeremy Aroles, and Wendelin Küpers extend the
enframing critique to educational technology.\footnote{Feldges and Pieczenko, ``Boredom in
Educational Contexts: A Critical Review,'' \textit{Encyclopaideia} 24.57 (2020); Kurt C. M. Mertel, ``Heidegger, Technology and Education,''
\textit{Journal of Philosophy of Education} 54.2 (2020); Aroles and K\"upers, ``Towards an Integral Pedagogy in the Age of `Digital
\textit{Gestell},''' 2022. The enframing argument these three extend is Iain Thomson's, whose
account of ontological education supplies the frame---\textit{paideia} against enframing,
teaching as letting learn, the university's crisis as ontohistorical rather than
institutional.}

Psychophysiology and affective computing supply the second literature's electroencephalographic
(EEG) work, and its most instructive results are those that failed. Joochan Kim, Jungryul Seo,
and Teemu Laine ran the closest precursor to this study's design: a monotonous video, a
consumer EEG headband and a screen-based eye tracker recording together, and a boredom rating
collected afterward. Their governing hypothesis was that the brain signal and the gaze signal
would agree with each other during boredom. Regions of temporally coincident change did appear
in both recordings across all thirteen analyzed participants, and the rating confirmed that
boredom had been elicited---but that rating correlated significantly with no extracted
physiological feature whatsoever. A study designed to demonstrate that the two recordings and the
person's own report converge instead documented their divergence, and that divergence is what
this section inherits and extends. Each such independent way of showing---what a person says
about an episode, what behavior displays during it, what an instrument records through it---is
a channel in what follows, and the point of Kim's result is that two channels trained on the
same episode can answer differently. Robert Barry and colleagues supply the methodological
floor beneath every EEG claim made here: eyes-closed and eyes-open rest are not interchangeable
baselines, since alpha power indexes global arousal while the focal bands index activation, so
a study must choose its baseline and defend the choice (2765--66, 2771--72). That bears on
this experiment twice, once on its EEG result and once on eye closure as a surface in its own
right. Rajamanickam Yuvaraj and colleagues are the nearest neighbor in both stimulus and
setting---university students, a boring lecture video, whole-scalp connectivity---and their
finding is that boredom is not a network-idle state: gamma-band connectivity grew more locally
clustered and less globally efficient, which they read as ``«Q16»'' rather than mere
disengagement (11). Their design also carries the gap this section must not repeat, since both
of their conditions are active eyes-open video-viewing with no resting baseline of any
kind.\footnote{Four further studies mark out what this work has and has not settled. Eri
Miyauchi and Masahiro Kawasaki timestamp the onset of boredom by having participants press a
key at the moment they first feel it, isolating becoming-bored from being-bored. Sammy Perone
and colleagues track frontal alpha asymmetry across a repetitive task and find it shifts as
the task wears on, making it a signature of ongoing emotional self-regulation rather than a
fixed trait marker. Ofir Yakobi and colleagues separate the routes by which state and trait
boredom show themselves, the boredom a task induces registering in task-evoked potentials and
the disposition to become bored only in resting oscillatory power. Jungryul Seo and colleagues classify boredom from a
two-electrode consumer headband at 86.7 percent accuracy, establishing that detection does not
require a dense montage.}

Eye tracking is the second measurement literature this study draws on, and it arrives with a
reporting standard, a warning about what the pupil actually indexes, and a warning about
blinks. Kenneth Holmqvist and forty-six co-authors set the standard: after reviewing what
instrument, environment, and participant factors demonstrably affect data quality, and
comparing that evidence against 207 published eye-tracking studies, they find reporting
practice governed largely by trust in the manufacturer's defaults, and derive a minimal
guideline of what a study must disclose for its numbers to be interpretable. This section's
apparatus and analysis are reported against that guideline. What the pupil indexes is the
second matter, and it is not boredom. Ruud van den Brink, Peter Murphy, and Sander Nieuwenhuis
tracked tonic pupil diameter against sustained-attention performance and found that attentional
lapses tended to occur when pupil diameter was small---but the finding did not survive in that
form. The apparent straight-line relation reflected the joint drift of pupil and performance
over time on task, and once that drift was statistically removed the relation became curved:
the most false alarms and the slowest, most variable responses occurred when the pupil was
both relatively small and relatively large, with performance best in between (9, 12). Too
little arousal and too much degrade attention alike. What the pupil indexes, then, is arousal
and attentional state, not what a person would say about being bored, which is why a pupil
result that fails to track self-reported boredom is an answer about the measure rather than a
failure of the experiment.
Mathias Scharinger, Yvonne Kammerer, and Peter Gerjets supply the sharper precedent, and it
concerns two measures rather than one.\footnote{Scharinger, Kammerer, and Gerjets, ``Pupil
Dilation and EEG Alpha Frequency Band Power Reveal Load on Executive Functions for
Link-Selection Processes during Text Reading'' (2015). The study asks what mental effort a
reader spends choosing which hyperlink to follow, and measures that effort two ways at once
across three experiments, the second and third built to rule out perceptual and motor
explanations of the first. It is not a study of boredom, and neither are Holmqvist's reporting
review or van den Brink's pupil work. All three enter this section as evidence about
measurement rather than about boredom itself: what an eye measure can be trusted to report,
and what follows when two measures of the same thing disagree.} Studying the load imposed by hyperlink selection during reading, they
recorded pupil dilation and EEG alpha power simultaneously. Both measures registered the load
increase as predicted. The two did not correlate with each other---$r$ = $-$.16 and $-$.14,
neither significant---and only alpha tracked reading comprehension. Two instruments can
therefore register the same manipulation, each validly, and still disagree with each other,
which makes the oppositional ranking this section reports a documented pattern in the
literature rather than an anomaly in this dataset. Gaze is the third measure and the one this
study leans on hardest, and the applied literature that uses it is almost uniformly
detection-oriented: Natasha Jaques and colleagues predict affective states from gaze features
during intelligent tutoring; Ruijie Wang, Yuanchen Xu, and Liming Chen's GazeMotive infers
motivation from gaze for students with learning difficulties; Kshitij Sharma, Michail
Giannakos, and Pierre Dillenbourg pair gaze with machine learning to raise motivation;
Saromporn Charoenpit and Michiko Ohkura read emotion from gaze in an e-learning system; Boris
Cheval and colleagues use it to measure attentional bias; and Vasileios Skaramagkas and
colleagues catalog the measures the whole enterprise runs on---how long and how often the eyes
fixate, how far and how fast they saccade, how widely gaze disperses across a scene, how often
and how long the lids close, and how the pupil changes---sorting each by whether it indexes
visual attention, emotional arousal, or mental workload, and concluding that no single measure
separates the three cleanly. What each of the applied studies builds from these measures is a
classifier: gaze features go in, an affective label comes out, and the label counts as correct
when it matches what the participant reported. That architecture installs self-report as the
arbiter of what gaze means, and it is the one commitment this study does not make. The
distinction is worth stating precisely, because the analysis reported below finds strong
agreement between gaze and self-report---close enough that a classifier built on this data
would likely work. That agreement is a measured result rather than an assumption, and keeping
it in the category of result is what allows the two to be read apart wherever they part
company. A design that has already installed self-report as the criterion cannot register a
divergence at all, on any film: it can only score the eye measure as having failed.
Blinks carry their own caution. Frans VanderWerf and colleagues measured spontaneous blinking
directly---a complete blink lasting 334 $\pm$ 67 ms, with the closing phase far quicker than
the reopening---and showed that duration and amplitude vary systematically with vertical gaze
position, so that where a person is looking changes what a blink measurement returns. Their
procedure records something more consequential still: the experimenters deliberately drew
participants' attention away from the experimental setting in order to sample blinks that were
genuinely spontaneous, which concedes that blinking is sensitive enough to attention that
attending to one's own situation alters it (2785--86). Any blink corpus recorded while a
stimulus is playing is consequently an unknown mixture of elicitation classes whose
composition depends on the stimulus itself.

Christine King and I designed this first experiment to identify and track boredom and
engagement by cross-validating three sources of evidence against one another: a survey
administered after each film, EEG recorded throughout, and eye tracking recorded throughout.
Beyond that stated aim, the study asked a further question---whether boredom, engagement,
attention, and immersion could be quantified at all, and whether the relations among them
could be measured rather than assumed. The expectation was that what a viewer reported, what
the brain signal showed, and what the eyes did would converge. Two of the three delivered. The eye measures separated the films and
tracked what viewers reported closely enough to be treated as findings in their own right. The
EEG did neither, though it did not come back empty: its most suggestive surface, default-mode
power, moved the same way in six of eight participants on both contrasts where the boring film
was involved, which is a possible relationship and is reported as one. It is not a result. No
EEG measure distinguishes the three films at conventional significance, and every pairwise
comparison falls well short once the arithmetic accounts for how many comparisons are being
made. The channel is also the study's thinnest, recorded in the first round only and never
repeated. What the EEG offers is a direction worth noting and nothing that can carry weight.
Set against the literature just surveyed, that outcome is characteristic rather than
exceptional. Kim's self-report correlated with no
physiological feature. Sammy
Perone and colleagues predicted that resting frontal asymmetry would differ according to which
kind of trait boredom a participant carried, and found no such relation at baseline.\footnote{
State boredom is the transient episode a situation occasions; trait boredom, or boredom
proneness, is an enduring disposition of the person to fall into it. Raffaelli and colleagues
scope their review of the empirical work to ``«Q29»'' (2452), the same division this section
observes: everything reported here concerns episodes occasioned by particular films, and no
claim is made about any participant's disposition.} Ofir Yakobi and colleagues expected boredom
to reveal itself twice over, in how well participants performed a sustained-attention task and
in the brain's response to each stimulus. Their account of why held that a bored person
withdraws attention rather than acting carelessly, so the errors should be failures to respond
to a target: the prediction was of ``«Q17».'' What they found instead was that the two kinds of
boredom took separate routes. The boredom the task induced showed up in the brain's responses
to stimuli. The standing disposition to become bored showed up only in resting brain activity,
recorded before the task began, and left task performance and the stimulus responses
untouched. The two never met. Craig Lindley and Lennart Nacke, inducing boredom,
immersion, and flow through the level design of a first-person shooter, found that the Game
Experience Questionnaire discriminated only two of its seven subscales, challenge and tension,
while the flow and immersion subscales---the very constructs the study was built
around---failed to separate the conditions the level design had demonstrably
produced.\footnote{The flow construct under test is Mihaly Csikszentmihalyi's, the ``«Q21»''
(qtd. in Lindley and Nacke 1--2).} Miyauchi and Kawasaki's oscillatory signature indexes the
appraisal a person makes of a task-unrelated thought rather than boredom itself. King and I
reached the same judgment about our own instrumentation and acted on it: after the pilot study
with eight participants we dropped EEG from the platform's assessment entirely, and our
reasons were those this literature keeps running into. On the signal itself we wrote that
``«Q26»,'' and the comparison with the alternative was decisive, since ``«Q27».'' Eye tracking
proved both the more accurate measure of whether a student was bored or engaged and the more
scalable one, because ``«Q28»'' (King and Salvo)---a laptop or tablet already has the camera,
where EEG needs an electrode array and a trained operator for every participant. That decision
was made on grounds of accuracy and scale, and made before the reprocessing reported in this
section had shown what the EEG measurements would and would not support.

% ★ BATCH 4 PARAGRAPHS DELETED BY AUTHOR RULING 2026-08-06. Both the lineage paragraph and
% the two-moves paragraph are cut: they do not bear on using Heideggerian phenomenology to
% analyze the physiological data, which is what this section does. Consequences recorded:
%   - Outline v7 §L1 ¶8 is DROPPED from the module. O-06 (how much to re-present vs cite) is
%     therefore MOOT for L1.
%   - The machine-indexed group entries no longer surface in prose: `grp-program-lineage`
%     (E15, extension-rather-than-resolution) and `grp-vr-phenomenology-framing` (E14, the
%     refs [13-15] gesture, invoked but inert). Q30 is now UNUSED.
%   - D-20's one-of-fifty-five observation is dropped with its host footnote.
%   - ★ CHECK AT L5: outline §B lists the desktop M7.1 hand-forward as REDEEMED partly by
%     "the unbuilt FCM-phenomenology x instrumentation combination (L1 supplies the warrant,
%     now definitional as well as bibliographic)". The BIBLIOGRAPHIC half survives in ¶9 and
%     the DEFINITIONAL half in ¶10, so the warrant still stands; only the program-lineage
%     narration is gone. No repair needed unless L5 wants the inert-framing point back.
%   - The study-design sentence the author asked for is folded into ¶7's opening instead.
%     (Q30 retained in the session bank though now unused; see errata.)
%   - The Fahlman definition now appears as a footnote at ¶1, cited to Fahlman (69) rather
%     than to the paper that adopted it, per author ruling.
%
% BATCH 5 = the two warrant paragraphs (outline v7 §L1 ¶¶9-10), drafted 2026-08-06.
%   ★ E02/E03 RE-PINNED against the regenerated synthesis layer this session, as the pack
%     required. Source: index/Boredom Secondary (Part III)/_synthesis/
%     boredom-construct-measure-concordance.md §2. Both routes CONFIRM UNCHANGED:
%       Route 1 (construct-based): 22 units treat profound boredom; 19 carry a hard
%         instrument (EEG / pupillometry / gaze / blink / ML); intersection EMPTY.
%       Route 2 (bridge-based): 21 units carry a bridges-to-fcm edge; 19 carry a
%         hard-instrument operationalizes-measure edge; intersection EMPTY.
%       ID ranges: FCM-bridging = bor-sec-01-18 (plus 23, 48, 51, 54, 55);
%         hard-instrumented = bor-sec-21-41. "Barely touch" is accurate.
%     Grid: 91 of 132 cells populated; the three FCM forms are reachable by SELF-REPORT ONLY
%       (every hard-instrument column empty across all three rows); the most important
%       populated cell (profound boredom x self-report, 5 units) is a NEGATIVE — bor-sec-04
%       argues the Boredom Proneness and Zuckerman scales do not measure profound boredom.
%   ¶10 renders the two definitional quotations TOGETHER AND ONCE per D-23: Eastwood's
%     condition (Q02, printed 484 — the body occurrence, not the barred abstract one) beside
%     the FCM's "nothing at all to be found" (Q03, printed 109). Q04 supplies the ontological
%     counter. Q22 is DISPLAYED per 0b. Q01 stays barred.
%   ★ FCM 109/110 RESOLVED — AUTHOR RULING 2026-08-07: the passage is on PAGE 109, and L1 is
%     correct as drafted. ★ DESKTOP DEBT: desktop M4.1 ¶6 cites 110 for the same passage and
%     must be corrected to 109 when that section is next opened.

% ★ BATCH 5 ¶8 (bibliographic warrant) DELETED BY AUTHOR RULING 2026-08-06: too metatextual,
% and the two fields and the gap between them are already established in ¶1. Outline v7 §L1 ¶9
% is therefore dropped as a paragraph; E02/E03's re-pinned counts are retained in this header
% for the works-cited and for L5 if the bibliographic half of the warrant is ever wanted.

Fifty-five works on boredom were surveyed for this section, and each was recorded twice over:
once for which conception of boredom it treats, and once for how, if at all, it measures
anything. Sorting them that way answers a question that decides what this experiment's data can
be asked. Of the fifty-five, nineteen measure with a physiological or computational
instrument---scalp electrodes, pupil diameter, gaze position, blink, or a trained
classifier---and every one of those nineteen studies a boredom that a person can notice, report,
and attribute to something. Not one of them studies the boredom Heidegger analyzes. The works
that do take up his analysis, twenty-two of them, measure nothing at all. The only way anyone
has reached Heidegger's forms empirically is by asking people, and asking is the one approach
that two separate bodies of work independently disqualify.

The phenomenological objection to asking is the more fundamental, and it is Heidegger's own. He
raises it about his own lecture course, which set out to awaken a fundamental attunement in his
listeners rather than to describe one. ``«Q32»'' (65), he writes, and then states the
prohibition:

\begin{quote}
«Q22». (65)
\end{quote}

\noindent The reason follows in his next sentence, and it is what makes the prohibition more
than a preference: ``«Q33»'' (65). To ascertain something is to bring it into awareness as an
object, to make it the thing one is now considering. An attunement, however, is not an object
within a situation but the way the whole situation is disclosed to the person in it, so
bringing it forward as an object changes what it was. That is the force of Heidegger's claim
that ``«Q23»'' (65). The alteration is ordinary enough to recognize: a person absorbed in an
afternoon, asked how they are finding it, must step out of the absorption to answer, and the
answer describes the stepping out as much as the absorption. Applied to this experiment the
consequence is exact. Each participant was asked, after each film, to rate boredom and
engagement on a seven-item survey. That survey is a legitimate measurement of what the
participant could recognize and report about the episode. It is not, and on Heidegger's account
cannot be, a measurement of an attunement, because the asking is itself an
intervention on what is being asked about.

The empirical objection to asking arrives from the opposite direction and needs no philosophy
at all. Kim, Seo, and Laine's participants rated their own boredom, and those ratings matched
none of the physiological measurements taken at the same time. Whatever the ratings recorded,
it was not what the electrodes and the eye tracker were recording. The survey is therefore the
only approach that reaches Heidegger's forms, and it is the approach that phenomenology says
cannot reach an attunement and that the measurement literature shows does not agree with the
physiological measures. It is worth being exact about what does and does not follow. The claim
is not that the surveyed works measured profound boredom badly. Nothing in the fifty-five
measured it. Where profound boredom and self-report appear together, they appear in works that
discuss both, and the most developed of those is Elpidorou and Freeman's argument that the
standard boredom scales, the Boredom Proneness Scale and the Boredom Susceptibility Scale among
them, do not measure Heidegger's boredom at all.

The reason these measures miss these forms is not technical, and it is worth following
carefully, because the reading this section performs depends on it. Every measure in that
survey was built to detect boredom as Eastwood defines it, and Eastwood's definition has three
conditions: a person is bored when they cannot engage attention in satisfying activity, when
they are aware of that failure, and, third, when they ``«Q02»'' (484). His own examples of the
third condition are ``this task is boring'' and ``there is nothing to do'' (484). Boredom, so
defined, always has an address. Something out there is responsible, and the bored person can
say what. Every measure built on that definition inherits the assumption, because each was
validated against people reporting exactly that kind of boredom.

Two of Heidegger's three forms withhold that address, and withhold it by construction rather
than by accident. As I articulated in the preceding section, the first form is boredom
\textit{by} something nameable: his scene is a four-hour wait at a small railway station, where
the station is what refuses, where the fight against the wait is visible in every glance at the
watch, and where the bored person could name the culprit if asked. That form satisfies
Eastwood's third condition exactly, and a measure calibrated on it will work. The second form
is different. Heidegger's scene is a dinner party where the food was good, the conversation
lively, the evening in no way unpleasant, and where only afterward does one realize one was
bored throughout. Asked at the table, the guest would have denied boredom and been telling the
truth. Nothing in the evening was responsible, and the guest could not attribute the state to
anything, precisely because there is ``«Q03»'' (109). Eastwood's third condition is not merely
unsatisfied here; it is refused on the definition's own terms, since no feature of the
surroundings is available to carry the cause. The third form withdraws further still.
Heidegger's phrase for it is that ``«Q04»'' (128): the boredom is occasioned by nothing in the
situation whatsoever, arising instead from the person's own being. A measure that asks what in
the surroundings is responsible has nothing to register in either case, not because the measure
is coarse, but because it puts a question these two forms give no answer to.

Two objections follow from all this, and they are different objections doing different work.
The first is about measurement. The only evidence that registers boredom as boredom is the
person's own statement that they were bored; nothing in a gaze trace or a band of scalp
activity announces itself as boredom, and every physiological pattern reported in this section
was identified by asking which measurements moved with those statements. The survey is
therefore the ground the rest of the analysis stands on, and the measurement question is what
follows from that dependence. The second objection is about definition, and it asks whether the
boredom those statements report---boredom as Eastwood's three conditions specify it, which is
what our survey items were written to elicit---is the boredom Heidegger analyzes at all. The
order matters, because the definitional question is settled before the measurement question can
be raised. If the forms this section reads were never within reach of a survey built on
environmental attribution, then no degree of agreement between that survey and any physiological
measurement bears on whether those forms were present. They were ruled out when the definition
was chosen, long before anything was recorded.

That dependence is worth stating precisely rather than apologizing for, because the analysis
below turns on how far it extends. Our original validation logic scored the physiological
measurements against what participants reported, and this section computes that agreement for
the first time. It is substantial. Gaze deviation tracks reported boredom at a rank correlation
of +0.590 and reported engagement at $-$0.585, both at Monte Carlo probabilities below 0.005;
eye closure follows at +0.460 and $-$0.416. The relationships are close to linear---the rank and
linear coefficients differ by a median of 0.032 across the whole family---so the more a
participant's gaze wandered from center, and the more of an episode their eyes spent closed, the
more bored that participant said they had been. Kim, Seo, and Laine, at thirteen participants,
found no such agreement with any measurement they took. At twelve participants here, several
associations hold. Where the earlier study could not connect what people said to what its
recordings showed, this one can.

What does not follow is that either kind of evidence certifies the other, and the pattern of
results shows why. Put two different questions to the same measurements and they answer in
different orders. Asked which measurement best matches what participants reported, gaze
deviation wins decisively. Asked which measurement best separates the three films from one
another, gaze deviation fails---its differences across films do not approach significance.
That combination is not a defect, and it is worth saying why, because the second question was
never this study's question. We did not set out to tell the three films apart by their
recordings; we set out to find what in the recordings moved with the felt response to them. A
measurement that tracks the felt response closely and sorts the films poorly is behaving
exactly as it should, because boredom is not a property of a film. It is what a particular
person makes of a particular seventeen minutes. S04 is the case that makes this concrete. She
rated the boring film 1 for boredom and 7 for engagement---the only participant to invert that
film so completely---while rating the clinical footage 9 and 1. Any measurement keyed to the
film would score her wrongly twice. A measurement keyed to what she reported follows her
inversion without difficulty, which is the whole point of scoring against reports rather than
against the schedule of what was shown.

Pupil dilation is the measurement that yields nothing this section builds on, and saying so
plainly is better than working around it. It does not track what participants reported---a rank
correlation of $-$0.172 against reported boredom, nowhere near significance---and while it does
distinguish the boring tutorial from the clinical footage, several uncontrolled factors could
produce that difference on their own. The samples were never matched for luminance, and pupils
respond to light before they respond to anything else. Mean dilation is in fact largest on the
clinical footage, at 3.51 mm against 3.40 for the interesting film and 3.32 for the boring
tutorial, which is the opposite of what the brightest sample should produce and suggests that
whatever is moving these pupils, it is not one thing. A measurement pulled by light, by arousal,
and by mental effort at once, and pulled by none of them decidedly here, is a measurement to
report and set aside. The argument that follows rests on the two measurements that did behave
consistently.

Eye closure is the measurement that answers both questions, and here too the separation is not
really about the films. It separates the boring tutorial from each of the other two at corrected
probabilities below 0.002, moving the same way for eleven of twelve participants against the
clinical footage and for all twelve against the interesting film; participants' eyes were closed
for roughly a quarter of the boring episodes on average and under a tenth of the others. It also
tracks what participants said, at +0.460 against reported boredom. Those two facts are one fact.
The boring tutorial drew the most closure because it drew the most boredom---mean reported
boredom of 6.9 against 5.4 for the clinical footage and 4.3 for the interesting film---so what
looks like a measurement sorting three videos is a measurement following the response those
videos happened to provoke in most people. Measurements that agreed on what boredom is would not
reorder like this depending on which question is asked. That they do is a finding about what each
is sensitive to, and it is reported as such rather than resolved by promoting one of them to the
truth of the matter.

The published corrective to the deficit reading of boredom comes from Elpidorou, who argues
that boredom does useful work: it informs a person that the present situation has stopped
satisfying, and it moves them toward one that might. ``«Q20»'' is his formulation, credited to
Fahlman and colleagues for the underlying claim (330).\footnote{Elpidorou's sentence carries
``(Fahlman et al., 2013, p. 68)''. The wording is his own---the word ``unstuck'' does not appear
in Fahlman's paper---and what he credits to Fahlman is the account of what boredom does, which
matters here because Fahlman is also the source of the definition our study adopted.} That
corrective is persuasive within its scope, and its scope is the whole of the point. It describes
a boredom that has an object, announces itself, and can be acted on, which is the first form.
Extending it to the second and third would be a mistake this section is careful not to make: a
boredom with no culprit gives nothing to move away from, and a boredom arising from one's own
being gives nowhere to move. The distinction drawn earlier between the episode a situation
occasions and the disposition a person carries guards the same ground from the other side.
Nothing reported here concerns anyone's disposition. Every measurement is of particular people
watching particular films on particular afternoons, and whether any of them is a person prone
to boredom is neither asked nor answerable from this data.
